\section{Rectangular boxes}\label{sec:boxes}
We prove Theorem~\ref{thm:all-boxes}. A prescribed-family corner lemma
and a uniform lower obstruction come first, followed by the thin-box
upper bounds and the remaining shapes. The only certified input is
Theorem~\ref{thm:main}; the prescribed-family lemma is independent of it.

After ordering the sides as $a\le b\le c$, the upper bound has three
exhaustive regimes:
\begin{center}
\begin{tabular}{ll}
\toprule
Shape & Main mechanism\\
\midrule
$a+b\ge6$ & Weighted compatibility and reservations\\
$a=b=2$ & Concentration and retargeting along a rod\\
$a=2,\ b=3$ & Exceptional-defect control and retargeting\\
\bottomrule
\end{tabular}
\end{center}
We establish the thin regimes first, then the general weighted bound.
In each case an initial corner matching verifies part of Hall's inequalities;
the final matching may rearrange those placements.

\begin{lemma}[Rare pairs]\label{lem:rare-pairs}
In a collection of five-colour cubes, any family of bichromatic pairs each
missing from more than half the source copies consists of at most two
pairwise disjoint pairs.
\end{lemma}
\begin{proof}
The missing families of two such pairs intersect. A source missing both
pairs forces them to be disjoint, by Lemma~\ref{lem:cube-incidence}.
There are at most two disjoint pairs in a five-colour palette.
\end{proof}

\subsection{Signed corner projections}
Retain the family $\mathcal H_{\alpha\beta}$ from Lemma~\ref{lem:fixed-missing-pair},
with its opposite-repeat varieties and their $2\times2$ corner array.
Write $C(W)$ for the cyclic corner types of $W$.

\begin{lemma}[Corner projections onto a missing-pair family]\label{lem:thin-projections}
Fix a bichromatic pair $\alpha\beta$.  Let $C^+$ and $C^-$ be the two sets of six
tricoloured corner types shared by the three adjacent-repeat varieties of
the respective cyclic orientations in $\mathcal H_{\alpha\beta}$.  An
opposite-repeat variety in $\mathcal H_{\alpha\beta}$ has four corners in each set.
For a source $W\notin\mathcal H_{\alpha\beta}$, put
\[
 p(W)=|C(W)\cap C^+|,\qquad m(W)=|C(W)\cap C^-|.
\]
Then
\[
 (p(W),m(W))\in\{(4,0),(2,1),(1,2),(2,2),(0,4)\}.
\]
\end{lemma}

\begin{proof}
Relabel the poles as $0,1$ and the ring colours as $2,3,4$.  Choose the
positive orientation so that
\[
 C^+=\{023,034,042,124,132,143\},\qquad
 C^-=\{024,032,043,123,134,142\}.
\]
Triples denote cyclic classes; face order is $(+x,-x,+y,-y,+z,-z)$.  The following table gives representatives of
all geometric cases outside $\mathcal H_{01}$, up to a colour permutation
preserving the pole pair and a simultaneous reflection, which may exchange
$C^+$ and $C^-$.  Such a reflection is only a symmetry of this incidence
calculation, not an allowed assembly move.
\[
\begin{array}{c|cc}
 W&p(W)&m(W)\\\hline
 (0,0,1,2,3,4)&2&2\\
 (0,1,0,2,3,4)&0&4\\
 (0,2,0,3,1,4)&2&1\\
 (0,2,1,2,3,4)&0&4\\
 (0,2,1,3,2,4)&2&1\\
 (0,2,1,3,4,4)&2&2
\end{array}
\]
To see that these cases exhaust the possibilities, first suppose a pole
colour is repeated.  Its two faces are either opposite, or adjacent; in
the latter case the other pole is opposite one repeated face or adjacent
to both.  If a ring colour is repeated, the two pole faces must be
adjacent, since otherwise $W$ belongs to $\mathcal H_{01}$.  Fix them at
$x^+$ and $y^+$.  Among the remaining faces $x^-,y^-,z^+,z^-$, an opposite
repeated pair is $z^+,z^-$.  An adjacent repeated pair is either $x^-,y^-$,
or one of $x^-,y^-$ together with one of $z^+,z^-$.  These are the final
three rows.

Each entry follows by reading the eight corners of the displayed cube.
The projections of the nine family
varieties themselves are $(6,0),(0,6)$, and $(4,4)$, as follows directly
from their ring descriptions.
\end{proof}

\begin{lemma}[Compatibility for a thin frame]\label{lem:thin-compatibility}
Choose an orientation of a five-colour target for a $2\times2\times3$
frame such that its four noncorner edge positions are bichromatic.
Every source cube fits at least four of the twelve positions.
\end{lemma}

\begin{proof}
If either source or target is opposite-repeat,
Lemma~\ref{lem:cube-incidence} gives two shared corners.
The four parallel bichromatic edges have either four distinct types or
two types occurring twice each. In the latter case both contain the
target's repeated colour. The source's missing pairs are disjoint, so
it fits at least two of these edges, giving four compatible positions.

If both varieties are adjacent-repeat, the four parallel edge types are
distinct and transverse to the monochromatic edge. The source fits either
all four, or at least three and a corner by
Lemma~\ref{lem:transverse-corner}. Again it fits four positions.
\end{proof}

\subsection{Eight family cubes and three arbitrary additions}

% Analytic lemma for the 2 x 2 x 3 five-colour problem.

\begin{lemma}[Eight cubes with a common missing pair]\label{lem:eight-common-missing}
Eight cubes from $\mathcal H_{\alpha\beta}$, together with any three additional
five-colour cubes, have a corner solution whose target belongs to
$\mathcal H_{\alpha\beta}$.
\end{lemma}

\begin{proof}
Let $O,P,Q$ count the opposite-repeat, positive
adjacent-repeat, and negative adjacent-repeat copies among the chosen eight.
Thus $O+P+Q=8$; $C^+$ and $C^-$ are as in Lemma~\ref{lem:thin-projections}.

If $O+\min(P,Q)\ge4$, the $2\times2$ Hall argument gives an
opposite-repeat corner solution from these eight copies alone. Explicitly,
choose $O_r$ whose opposite-repeat multiplicity is maximal. Each cell has
at least $O/2+\min(P,Q)\ge2$ neighbours, each row has at least
$8-O/2\ge4$, and each column has at least $O+\min(P,Q)\ge4$.
Every other nonempty union of cells has all eight copies as neighbours.

Otherwise exchange the signs so that $O+Q\le3$. Put
$\kappa=O+Q=8-P$, so $P\ge5$ and $0\le \kappa\le3$.
Some positive variety $A_r^+$ occurs at least twice. If its eight
corners have a matching among all eleven copies, we are done. Otherwise,
choose a deficient Hall set $D$ of its corners. The two copies of
$A_r^+$ fit every corner, so $|D|\ge3$ and $D$ contains a tricoloured
corner. All $P$ positive copies are therefore neighbours of $D$.
Let $R$ be its other neighbours and put $U=C^+\setminus D$, $u=|U|$.
Since $|D|+u\le8$,
\[
 P+|R|<|D|\le8-u,\qquad |R|+u\le \kappa-1.
\]
Remove the $P$ positive copies and the copies in $R$. Every surviving
copy has all its positive corners in $U$, and $u\le2$.

Every family opposite-repeat variety has four corners in $C^+$, so no
such chosen copy remains. At least
\[
 \kappa-|R|\ge u+1
\]
of the chosen nonpositive copies remain, and all of them are negative
adjacent-repeat copies. Each fits every corner in $C^-$. The remaining
pool has at least
\[
 \kappa+3-|R|\ge4+u
\]
copies.

Lemma~\ref{lem:thin-projections} gives the following implications for every source:
\[
\begin{array}{c|ccc}
 |C(W)\cap C^+|\le u &u=0&u=1&u=2\\\hline
 |C(W)\cap C^-|\ge&4&2&1.
\end{array}
\]
Indeed, the family varieties have projection sizes $(6,0),(0,6),(4,4)$,
and every variety outside the family has one of
$(4,0),(2,1),(1,2),(2,2),(0,4)$.

Partition $C^-$ into three pairs, indexed by the unordered pairs of ring
colours. The negative half of each family opposite-repeat target is the
union of two of these pairs. We will match one such four-position half
using the remaining pool, then use four positive copies for its positive
half.

Let $j$ be the number of negative-family copies in the remaining pool;
thus $j\ge u+1$. If $j\ge4$, the conclusion is immediate. If $j=3$,
another remaining copy fits at least one negative corner; choose a target
whose negative half contains it and fill the other three positions with
negative-family copies.

If $j=2$, then $u\le1$. Two other remaining copies each fit at least two
negative corners, so they can be assigned distinct ones. Any two negative
corners lie in the union of two of the three pairs. Choose the corresponding
opposite-repeat target and fill its other two negative positions with the
two negative-family copies.

Finally, if $j=1$, then $u=0$, and there are at least three other copies,
each fitting at least four of the six negative corners. Choose one of
these copies $W$. Some union of two of the three pairs contains at least
three of its negative corners. In that four-position half, every other
copy fits at least two positions. Hence $W$ and any two other copies can
be matched to three distinct positions: two copies have at least two
neighbours, and all three together have at least three. The one
negative-family copy fills the fourth position.

In every case we obtain the negative half of an opposite-repeat target
in $\mathcal H_{\alpha\beta}$. Four of the $P\ge5$ positive copies fill its
positive half. The positive copies used here are disjoint from the surviving pool. This proves the lemma.
\end{proof}


\subsection{A uniform lower bound}

\begin{lemma}[A three-variety obstruction]\label{lem:thin-three-variety-obstruction}
Let $a,b,c\ge2$ and put $\tau=a+b+c-6$. In face order
$(+x,-x,+y,-y,+z,-z)$, define
\[
 U=(0,1,0,2,3,4),\qquad
 V=(0,1,1,2,3,4),\qquad
 W=(0,3,1,4,2,3).
\]
The collection
\[
 U+(\tau+5)V+(\tau+6)W
\]
of $2\tau+12$ cubes has no $a\times b\times c$ frame using all five
colours. In particular, $U+7V+8W$ fails for both a $2\times3\times3$
frame and a $2\times2\times4$ box.
\end{lemma}

\begin{proof}
The corner types are
\[
\begin{array}{c|l}
 U&003,004,013,024,032,041,123,142\\
 V&113,114,013,024,032,041,123,142\\
 W&133,334,012,024,031,043,132,234.
\end{array}
\]
Triples denote cyclic classes. The only two-colour types in their union
are $003,004$, $113,114$, and $133,334$. An adjacent-repeat target must
therefore repeat colour $0$, $1$, or $3$. Its two repeated-colour corners
force the colours on the opposite faces at the ends of its monochromatic
edge to be, respectively, $\{3,4\}$, $\{3,4\}$, or $\{1,4\}$.
The two orders of the remaining colours give $U,V,W$, respectively, and
their mirrors. The mirrors of $U$ and $V$ require $023$, which is absent
from the union; the mirror of $W$ requires $021$, also absent.

An opposite-repeat target requires both cyclic orientations of four
distinct unordered tricoloured triples. The union above contains both
orientations only for $013$ and $123$. Thus no opposite-repeat target
is possible. The only corner-supported targets are $U,V,W$.

Target $U$ needs two copies possessing a $00$ corner, but the collection
has only one such copy.

For target $V$, seven of its corner types are absent from $W$; the only
shared type is $024$. In addition, exactly one geometric edge in each
coordinate direction has a type absent from $W$: these are $14,14,11$.
Consequently, in any orientation of an $a\times b\times c$ frame,
$7+\tau$ positions require a $U$ or $V$ source. There are only $\tau+6$
such copies.

For target $W$, seven corner types are absent from both $U$ and $V$.
Again, one geometric edge in each direction has a private type, now
$34,33,34$. Hence $7+\tau$ positions require a $W$ source, whereas only
$\tau+6$ are available.

Every possible target therefore fails
(Figure~\ref{fig:uniform-obstruction}).
\end{proof}

\begin{figure}[ht]
\centering
% Generated by scripts/generate_five_colour_figure.py; do not edit.
\definecolor{cube0}{HTML}{E69F00}
\definecolor{cube1}{HTML}{56B4E9}
\definecolor{cube2}{HTML}{009E73}
\definecolor{cube3}{HTML}{F0E442}
\definecolor{cube4}{HTML}{CC79A7}
\begin{tikzpicture}[x=0.176mm,y=-0.176mm]
\path[use as bounding box] (0,0) rectangle (900,470);
\node[anchor=center,inner sep=0pt,outer sep=0pt,text=black!93,font={\sffamily\fontsize{8}{9.6}\selectfont\bfseries}] at (72,22) {Collection};
\node[anchor=center,inner sep=0pt,outer sep=0pt,text=black!93,font={\sffamily\fontsize{8}{9.6}\selectfont\bfseries}] at (348,22) {Corner types of each target};
\node[anchor=center,inner sep=0pt,outer sep=0pt,text=black!93,font={\sffamily\fontsize{8}{9.6}\selectfont\bfseries}] at (624,22) {Edge types};
\node[anchor=west,inner sep=0pt,outer sep=0pt,text=black!93,font={\sffamily\fontsize{8}{9.6}\selectfont\bfseries}] at (716,22) {Why it fails};
\draw[draw=black!20,line width=0.599pt] (16,40) -- (884,40);
\draw[fill=cube3,draw=black!87,line join=round,line width=0.599pt] (48,54) -- (72,54) -- (72,78) -- (48,78) -- cycle;
\node[anchor=center,inner sep=0pt,outer sep=0pt,text=black!93,font={\sffamily\fontsize{8}{9.6}\selectfont}] at (60,66) {3};
\draw[fill=cube2,draw=black!87,line join=round,line width=0.599pt] (24,78) -- (48,78) -- (48,102) -- (24,102) -- cycle;
\node[anchor=center,inner sep=0pt,outer sep=0pt,text=black!93,font={\sffamily\fontsize{8}{9.6}\selectfont}] at (36,90) {2};
\draw[fill=cube0,draw=black!87,line join=round,line width=0.599pt] (48,78) -- (72,78) -- (72,102) -- (48,102) -- cycle;
\node[anchor=center,inner sep=0pt,outer sep=0pt,text=black!93,font={\sffamily\fontsize{8}{9.6}\selectfont}] at (60,90) {0};
\draw[fill=cube0,draw=black!87,line join=round,line width=0.599pt] (72,78) -- (96,78) -- (96,102) -- (72,102) -- cycle;
\node[anchor=center,inner sep=0pt,outer sep=0pt,text=black!93,font={\sffamily\fontsize{8}{9.6}\selectfont}] at (84,90) {0};
\draw[fill=cube1,draw=black!87,line join=round,line width=0.599pt] (96,78) -- (120,78) -- (120,102) -- (96,102) -- cycle;
\node[anchor=center,inner sep=0pt,outer sep=0pt,text=black!93,font={\sffamily\fontsize{8}{9.6}\selectfont}] at (108,90) {1};
\draw[fill=cube4,draw=black!87,line join=round,line width=0.599pt] (48,102) -- (72,102) -- (72,126) -- (48,126) -- cycle;
\node[anchor=center,inner sep=0pt,outer sep=0pt,text=black!93,font={\sffamily\fontsize{8}{9.6}\selectfont}] at (60,114) {4};
\node[anchor=center,inner sep=0pt,outer sep=0pt,text=black!93,font={\sffamily\fontsize{9}{10.8}\selectfont\bfseries}] at (72,148) {1 \(\times\) U};
\node[anchor=center,inner sep=0pt,outer sep=0pt,text=black!93,font={\sffamily\fontsize{8}{9.6}\selectfont}] at (72,170) {010234};
\draw[fill=cube0,draw=black!87,line join=round,line width=0.599pt] (180,96) -- (162.679,86) -- (162.679,106) -- (180,116) -- cycle;
\node[anchor=center,inner sep=0pt,outer sep=0pt,text=black!93,font={\sffamily\fontsize{8}{9.6}\selectfont}] at (171.34,101) {0};
\draw[fill=cube0,draw=black!87,line join=round,line width=0.599pt] (180,96) -- (180,116) -- (197.321,106) -- (197.321,86) -- cycle;
\node[anchor=center,inner sep=0pt,outer sep=0pt,text=black!93,font={\sffamily\fontsize{8}{9.6}\selectfont}] at (188.66,101) {0};
\draw[fill=cube3,draw=black!87,line join=round,line width=0.599pt] (180,96) -- (197.321,86) -- (180,76) -- (162.679,86) -- cycle;
\node[anchor=center,inner sep=0pt,outer sep=0pt,text=black!93,font={\sffamily\fontsize{8}{9.6}\selectfont}] at (180,86) {3};
\node[anchor=center,inner sep=0pt,outer sep=0pt,text=black!93,font={\sffamily\fontsize{8}{9.6}\selectfont}] at (180,131) {003};
\draw[fill=cube0,draw=black!87,line join=round,line width=0.599pt] (228,96) -- (210.679,86) -- (210.679,106) -- (228,116) -- cycle;
\node[anchor=center,inner sep=0pt,outer sep=0pt,text=black!93,font={\sffamily\fontsize{8}{9.6}\selectfont}] at (219.34,101) {0};
\draw[fill=cube0,draw=black!87,line join=round,line width=0.599pt] (228,96) -- (228,116) -- (245.321,106) -- (245.321,86) -- cycle;
\node[anchor=center,inner sep=0pt,outer sep=0pt,text=black!93,font={\sffamily\fontsize{8}{9.6}\selectfont}] at (236.66,101) {0};
\draw[fill=cube4,draw=black!87,line join=round,line width=0.599pt] (228,96) -- (245.321,86) -- (228,76) -- (210.679,86) -- cycle;
\node[anchor=center,inner sep=0pt,outer sep=0pt,text=black!93,font={\sffamily\fontsize{8}{9.6}\selectfont}] at (228,86) {4};
\node[anchor=center,inner sep=0pt,outer sep=0pt,text=black!93,font={\sffamily\fontsize{8}{9.6}\selectfont}] at (228,131) {004};
\draw[fill=cube0!30,draw=black!40,line join=round,line width=0.599pt] (276,96) -- (258.679,86) -- (258.679,106) -- (276,116) -- cycle;
\node[anchor=center,inner sep=0pt,outer sep=0pt,text=black!60,font={\sffamily\fontsize{8}{9.6}\selectfont}] at (267.34,101) {0};
\draw[fill=cube1!30,draw=black!40,line join=round,line width=0.599pt] (276,96) -- (276,116) -- (293.321,106) -- (293.321,86) -- cycle;
\node[anchor=center,inner sep=0pt,outer sep=0pt,text=black!60,font={\sffamily\fontsize{8}{9.6}\selectfont}] at (284.66,101) {1};
\draw[fill=cube3!30,draw=black!40,line join=round,line width=0.599pt] (276,96) -- (293.321,86) -- (276,76) -- (258.679,86) -- cycle;
\node[anchor=center,inner sep=0pt,outer sep=0pt,text=black!60,font={\sffamily\fontsize{8}{9.6}\selectfont}] at (276,86) {3};
\node[anchor=center,inner sep=0pt,outer sep=0pt,text=black!60,font={\sffamily\fontsize{8}{9.6}\selectfont}] at (276,131) {013};
\draw[fill=cube0!30,draw=black!40,line join=round,line width=0.599pt] (324,96) -- (306.679,86) -- (306.679,106) -- (324,116) -- cycle;
\node[anchor=center,inner sep=0pt,outer sep=0pt,text=black!60,font={\sffamily\fontsize{8}{9.6}\selectfont}] at (315.34,101) {0};
\draw[fill=cube2!30,draw=black!40,line join=round,line width=0.599pt] (324,96) -- (324,116) -- (341.321,106) -- (341.321,86) -- cycle;
\node[anchor=center,inner sep=0pt,outer sep=0pt,text=black!60,font={\sffamily\fontsize{8}{9.6}\selectfont}] at (332.66,101) {2};
\draw[fill=cube4!30,draw=black!40,line join=round,line width=0.599pt] (324,96) -- (341.321,86) -- (324,76) -- (306.679,86) -- cycle;
\node[anchor=center,inner sep=0pt,outer sep=0pt,text=black!60,font={\sffamily\fontsize{8}{9.6}\selectfont}] at (324,86) {4};
\node[anchor=center,inner sep=0pt,outer sep=0pt,text=black!60,font={\sffamily\fontsize{8}{9.6}\selectfont}] at (324,131) {024};
\draw[fill=cube0!30,draw=black!40,line join=round,line width=0.599pt] (372,96) -- (354.679,86) -- (354.679,106) -- (372,116) -- cycle;
\node[anchor=center,inner sep=0pt,outer sep=0pt,text=black!60,font={\sffamily\fontsize{8}{9.6}\selectfont}] at (363.34,101) {0};
\draw[fill=cube3!30,draw=black!40,line join=round,line width=0.599pt] (372,96) -- (372,116) -- (389.321,106) -- (389.321,86) -- cycle;
\node[anchor=center,inner sep=0pt,outer sep=0pt,text=black!60,font={\sffamily\fontsize{8}{9.6}\selectfont}] at (380.66,101) {3};
\draw[fill=cube2!30,draw=black!40,line join=round,line width=0.599pt] (372,96) -- (389.321,86) -- (372,76) -- (354.679,86) -- cycle;
\node[anchor=center,inner sep=0pt,outer sep=0pt,text=black!60,font={\sffamily\fontsize{8}{9.6}\selectfont}] at (372,86) {2};
\node[anchor=center,inner sep=0pt,outer sep=0pt,text=black!60,font={\sffamily\fontsize{8}{9.6}\selectfont}] at (372,131) {032};
\draw[fill=cube0!30,draw=black!40,line join=round,line width=0.599pt] (420,96) -- (402.679,86) -- (402.679,106) -- (420,116) -- cycle;
\node[anchor=center,inner sep=0pt,outer sep=0pt,text=black!60,font={\sffamily\fontsize{8}{9.6}\selectfont}] at (411.34,101) {0};
\draw[fill=cube4!30,draw=black!40,line join=round,line width=0.599pt] (420,96) -- (420,116) -- (437.321,106) -- (437.321,86) -- cycle;
\node[anchor=center,inner sep=0pt,outer sep=0pt,text=black!60,font={\sffamily\fontsize{8}{9.6}\selectfont}] at (428.66,101) {4};
\draw[fill=cube1!30,draw=black!40,line join=round,line width=0.599pt] (420,96) -- (437.321,86) -- (420,76) -- (402.679,86) -- cycle;
\node[anchor=center,inner sep=0pt,outer sep=0pt,text=black!60,font={\sffamily\fontsize{8}{9.6}\selectfont}] at (420,86) {1};
\node[anchor=center,inner sep=0pt,outer sep=0pt,text=black!60,font={\sffamily\fontsize{8}{9.6}\selectfont}] at (420,131) {041};
\draw[fill=cube1!30,draw=black!40,line join=round,line width=0.599pt] (468,96) -- (450.679,86) -- (450.679,106) -- (468,116) -- cycle;
\node[anchor=center,inner sep=0pt,outer sep=0pt,text=black!60,font={\sffamily\fontsize{8}{9.6}\selectfont}] at (459.34,101) {1};
\draw[fill=cube2!30,draw=black!40,line join=round,line width=0.599pt] (468,96) -- (468,116) -- (485.321,106) -- (485.321,86) -- cycle;
\node[anchor=center,inner sep=0pt,outer sep=0pt,text=black!60,font={\sffamily\fontsize{8}{9.6}\selectfont}] at (476.66,101) {2};
\draw[fill=cube3!30,draw=black!40,line join=round,line width=0.599pt] (468,96) -- (485.321,86) -- (468,76) -- (450.679,86) -- cycle;
\node[anchor=center,inner sep=0pt,outer sep=0pt,text=black!60,font={\sffamily\fontsize{8}{9.6}\selectfont}] at (468,86) {3};
\node[anchor=center,inner sep=0pt,outer sep=0pt,text=black!60,font={\sffamily\fontsize{8}{9.6}\selectfont}] at (468,131) {123};
\draw[fill=cube1!30,draw=black!40,line join=round,line width=0.599pt] (516,96) -- (498.679,86) -- (498.679,106) -- (516,116) -- cycle;
\node[anchor=center,inner sep=0pt,outer sep=0pt,text=black!60,font={\sffamily\fontsize{8}{9.6}\selectfont}] at (507.34,101) {1};
\draw[fill=cube4!30,draw=black!40,line join=round,line width=0.599pt] (516,96) -- (516,116) -- (533.321,106) -- (533.321,86) -- cycle;
\node[anchor=center,inner sep=0pt,outer sep=0pt,text=black!60,font={\sffamily\fontsize{8}{9.6}\selectfont}] at (524.66,101) {4};
\draw[fill=cube2!30,draw=black!40,line join=round,line width=0.599pt] (516,96) -- (533.321,86) -- (516,76) -- (498.679,86) -- cycle;
\node[anchor=center,inner sep=0pt,outer sep=0pt,text=black!60,font={\sffamily\fontsize{8}{9.6}\selectfont}] at (516,86) {2};
\node[anchor=center,inner sep=0pt,outer sep=0pt,text=black!60,font={\sffamily\fontsize{8}{9.6}\selectfont}] at (516,131) {142};
\draw[draw=black!87,line width=0.699pt] (158,144) -- (158,150) -- (250,150) -- (250,144);
\node[anchor=center,inner sep=0pt,outer sep=0pt,text=black!93,font={\sffamily\fontsize{8}{9.6}\selectfont\bfseries}] at (204,168) {2 private};
\node[anchor=west,inner sep=0pt,outer sep=0pt,text=black!93,font={\sffamily\fontsize{8}{9.6}\selectfont}] at (716,76) {2 corners};
\node[anchor=west,inner sep=0pt,outer sep=0pt,text=black!93,font={\sffamily\fontsize{8}{9.6}\selectfont}] at (716,100) {need U;};
\node[anchor=west,inner sep=0pt,outer sep=0pt,text=black!93,font={\sffamily\fontsize{8}{9.6}\selectfont\bfseries}] at (716,124) {only 1 copy.};
\draw[draw=black!20,line width=0.599pt] (16,180) -- (884,180);
\draw[fill=cube3,draw=black!87,line join=round,line width=0.599pt] (48,194) -- (72,194) -- (72,218) -- (48,218) -- cycle;
\node[anchor=center,inner sep=0pt,outer sep=0pt,text=black!93,font={\sffamily\fontsize{8}{9.6}\selectfont}] at (60,206) {3};
\draw[fill=cube2,draw=black!87,line join=round,line width=0.599pt] (24,218) -- (48,218) -- (48,242) -- (24,242) -- cycle;
\node[anchor=center,inner sep=0pt,outer sep=0pt,text=black!93,font={\sffamily\fontsize{8}{9.6}\selectfont}] at (36,230) {2};
\draw[fill=cube0,draw=black!87,line join=round,line width=0.599pt] (48,218) -- (72,218) -- (72,242) -- (48,242) -- cycle;
\node[anchor=center,inner sep=0pt,outer sep=0pt,text=black!93,font={\sffamily\fontsize{8}{9.6}\selectfont}] at (60,230) {0};
\draw[fill=cube1,draw=black!87,line join=round,line width=0.599pt] (72,218) -- (96,218) -- (96,242) -- (72,242) -- cycle;
\node[anchor=center,inner sep=0pt,outer sep=0pt,text=black!93,font={\sffamily\fontsize{8}{9.6}\selectfont}] at (84,230) {1};
\draw[fill=cube1,draw=black!87,line join=round,line width=0.599pt] (96,218) -- (120,218) -- (120,242) -- (96,242) -- cycle;
\node[anchor=center,inner sep=0pt,outer sep=0pt,text=black!93,font={\sffamily\fontsize{8}{9.6}\selectfont}] at (108,230) {1};
\draw[fill=cube4,draw=black!87,line join=round,line width=0.599pt] (48,242) -- (72,242) -- (72,266) -- (48,266) -- cycle;
\node[anchor=center,inner sep=0pt,outer sep=0pt,text=black!93,font={\sffamily\fontsize{8}{9.6}\selectfont}] at (60,254) {4};
\node[anchor=center,inner sep=0pt,outer sep=0pt,text=black!93,font={\sffamily\fontsize{9}{10.8}\selectfont\bfseries}] at (72,288) {(\(\tau\)+5) \(\times\) V};
\node[anchor=center,inner sep=0pt,outer sep=0pt,text=black!93,font={\sffamily\fontsize{8}{9.6}\selectfont}] at (72,310) {011234};
\draw[fill=cube0,draw=black!87,line join=round,line width=0.599pt] (180,236) -- (162.679,226) -- (162.679,246) -- (180,256) -- cycle;
\node[anchor=center,inner sep=0pt,outer sep=0pt,text=black!93,font={\sffamily\fontsize{8}{9.6}\selectfont}] at (171.34,241) {0};
\draw[fill=cube1,draw=black!87,line join=round,line width=0.599pt] (180,236) -- (180,256) -- (197.321,246) -- (197.321,226) -- cycle;
\node[anchor=center,inner sep=0pt,outer sep=0pt,text=black!93,font={\sffamily\fontsize{8}{9.6}\selectfont}] at (188.66,241) {1};
\draw[fill=cube3,draw=black!87,line join=round,line width=0.599pt] (180,236) -- (197.321,226) -- (180,216) -- (162.679,226) -- cycle;
\node[anchor=center,inner sep=0pt,outer sep=0pt,text=black!93,font={\sffamily\fontsize{8}{9.6}\selectfont}] at (180,226) {3};
\node[anchor=center,inner sep=0pt,outer sep=0pt,text=black!93,font={\sffamily\fontsize{8}{9.6}\selectfont}] at (180,271) {013};
\draw[fill=cube0,draw=black!87,line join=round,line width=0.599pt] (228,236) -- (210.679,226) -- (210.679,246) -- (228,256) -- cycle;
\node[anchor=center,inner sep=0pt,outer sep=0pt,text=black!93,font={\sffamily\fontsize{8}{9.6}\selectfont}] at (219.34,241) {0};
\draw[fill=cube3,draw=black!87,line join=round,line width=0.599pt] (228,236) -- (228,256) -- (245.321,246) -- (245.321,226) -- cycle;
\node[anchor=center,inner sep=0pt,outer sep=0pt,text=black!93,font={\sffamily\fontsize{8}{9.6}\selectfont}] at (236.66,241) {3};
\draw[fill=cube2,draw=black!87,line join=round,line width=0.599pt] (228,236) -- (245.321,226) -- (228,216) -- (210.679,226) -- cycle;
\node[anchor=center,inner sep=0pt,outer sep=0pt,text=black!93,font={\sffamily\fontsize{8}{9.6}\selectfont}] at (228,226) {2};
\node[anchor=center,inner sep=0pt,outer sep=0pt,text=black!93,font={\sffamily\fontsize{8}{9.6}\selectfont}] at (228,271) {032};
\draw[fill=cube0,draw=black!87,line join=round,line width=0.599pt] (276,236) -- (258.679,226) -- (258.679,246) -- (276,256) -- cycle;
\node[anchor=center,inner sep=0pt,outer sep=0pt,text=black!93,font={\sffamily\fontsize{8}{9.6}\selectfont}] at (267.34,241) {0};
\draw[fill=cube4,draw=black!87,line join=round,line width=0.599pt] (276,236) -- (276,256) -- (293.321,246) -- (293.321,226) -- cycle;
\node[anchor=center,inner sep=0pt,outer sep=0pt,text=black!93,font={\sffamily\fontsize{8}{9.6}\selectfont}] at (284.66,241) {4};
\draw[fill=cube1,draw=black!87,line join=round,line width=0.599pt] (276,236) -- (293.321,226) -- (276,216) -- (258.679,226) -- cycle;
\node[anchor=center,inner sep=0pt,outer sep=0pt,text=black!93,font={\sffamily\fontsize{8}{9.6}\selectfont}] at (276,226) {1};
\node[anchor=center,inner sep=0pt,outer sep=0pt,text=black!93,font={\sffamily\fontsize{8}{9.6}\selectfont}] at (276,271) {041};
\draw[fill=cube1,draw=black!87,line join=round,line width=0.599pt] (324,236) -- (306.679,226) -- (306.679,246) -- (324,256) -- cycle;
\node[anchor=center,inner sep=0pt,outer sep=0pt,text=black!93,font={\sffamily\fontsize{8}{9.6}\selectfont}] at (315.34,241) {1};
\draw[fill=cube1,draw=black!87,line join=round,line width=0.599pt] (324,236) -- (324,256) -- (341.321,246) -- (341.321,226) -- cycle;
\node[anchor=center,inner sep=0pt,outer sep=0pt,text=black!93,font={\sffamily\fontsize{8}{9.6}\selectfont}] at (332.66,241) {1};
\draw[fill=cube3,draw=black!87,line join=round,line width=0.599pt] (324,236) -- (341.321,226) -- (324,216) -- (306.679,226) -- cycle;
\node[anchor=center,inner sep=0pt,outer sep=0pt,text=black!93,font={\sffamily\fontsize{8}{9.6}\selectfont}] at (324,226) {3};
\node[anchor=center,inner sep=0pt,outer sep=0pt,text=black!93,font={\sffamily\fontsize{8}{9.6}\selectfont}] at (324,271) {113};
\draw[fill=cube1,draw=black!87,line join=round,line width=0.599pt] (372,236) -- (354.679,226) -- (354.679,246) -- (372,256) -- cycle;
\node[anchor=center,inner sep=0pt,outer sep=0pt,text=black!93,font={\sffamily\fontsize{8}{9.6}\selectfont}] at (363.34,241) {1};
\draw[fill=cube1,draw=black!87,line join=round,line width=0.599pt] (372,236) -- (372,256) -- (389.321,246) -- (389.321,226) -- cycle;
\node[anchor=center,inner sep=0pt,outer sep=0pt,text=black!93,font={\sffamily\fontsize{8}{9.6}\selectfont}] at (380.66,241) {1};
\draw[fill=cube4,draw=black!87,line join=round,line width=0.599pt] (372,236) -- (389.321,226) -- (372,216) -- (354.679,226) -- cycle;
\node[anchor=center,inner sep=0pt,outer sep=0pt,text=black!93,font={\sffamily\fontsize{8}{9.6}\selectfont}] at (372,226) {4};
\node[anchor=center,inner sep=0pt,outer sep=0pt,text=black!93,font={\sffamily\fontsize{8}{9.6}\selectfont}] at (372,271) {114};
\draw[fill=cube1,draw=black!87,line join=round,line width=0.599pt] (420,236) -- (402.679,226) -- (402.679,246) -- (420,256) -- cycle;
\node[anchor=center,inner sep=0pt,outer sep=0pt,text=black!93,font={\sffamily\fontsize{8}{9.6}\selectfont}] at (411.34,241) {1};
\draw[fill=cube2,draw=black!87,line join=round,line width=0.599pt] (420,236) -- (420,256) -- (437.321,246) -- (437.321,226) -- cycle;
\node[anchor=center,inner sep=0pt,outer sep=0pt,text=black!93,font={\sffamily\fontsize{8}{9.6}\selectfont}] at (428.66,241) {2};
\draw[fill=cube3,draw=black!87,line join=round,line width=0.599pt] (420,236) -- (437.321,226) -- (420,216) -- (402.679,226) -- cycle;
\node[anchor=center,inner sep=0pt,outer sep=0pt,text=black!93,font={\sffamily\fontsize{8}{9.6}\selectfont}] at (420,226) {3};
\node[anchor=center,inner sep=0pt,outer sep=0pt,text=black!93,font={\sffamily\fontsize{8}{9.6}\selectfont}] at (420,271) {123};
\draw[fill=cube1,draw=black!87,line join=round,line width=0.599pt] (468,236) -- (450.679,226) -- (450.679,246) -- (468,256) -- cycle;
\node[anchor=center,inner sep=0pt,outer sep=0pt,text=black!93,font={\sffamily\fontsize{8}{9.6}\selectfont}] at (459.34,241) {1};
\draw[fill=cube4,draw=black!87,line join=round,line width=0.599pt] (468,236) -- (468,256) -- (485.321,246) -- (485.321,226) -- cycle;
\node[anchor=center,inner sep=0pt,outer sep=0pt,text=black!93,font={\sffamily\fontsize{8}{9.6}\selectfont}] at (476.66,241) {4};
\draw[fill=cube2,draw=black!87,line join=round,line width=0.599pt] (468,236) -- (485.321,226) -- (468,216) -- (450.679,226) -- cycle;
\node[anchor=center,inner sep=0pt,outer sep=0pt,text=black!93,font={\sffamily\fontsize{8}{9.6}\selectfont}] at (468,226) {2};
\node[anchor=center,inner sep=0pt,outer sep=0pt,text=black!93,font={\sffamily\fontsize{8}{9.6}\selectfont}] at (468,271) {142};
\draw[fill=cube0!30,draw=black!40,line join=round,line width=0.599pt] (516,236) -- (498.679,226) -- (498.679,246) -- (516,256) -- cycle;
\node[anchor=center,inner sep=0pt,outer sep=0pt,text=black!60,font={\sffamily\fontsize{8}{9.6}\selectfont}] at (507.34,241) {0};
\draw[fill=cube2!30,draw=black!40,line join=round,line width=0.599pt] (516,236) -- (516,256) -- (533.321,246) -- (533.321,226) -- cycle;
\node[anchor=center,inner sep=0pt,outer sep=0pt,text=black!60,font={\sffamily\fontsize{8}{9.6}\selectfont}] at (524.66,241) {2};
\draw[fill=cube4!30,draw=black!40,line join=round,line width=0.599pt] (516,236) -- (533.321,226) -- (516,216) -- (498.679,226) -- cycle;
\node[anchor=center,inner sep=0pt,outer sep=0pt,text=black!60,font={\sffamily\fontsize{8}{9.6}\selectfont}] at (516,226) {4};
\node[anchor=center,inner sep=0pt,outer sep=0pt,text=black!60,font={\sffamily\fontsize{8}{9.6}\selectfont}] at (516,271) {024};
\draw[draw=black!87,line width=0.699pt] (158,284) -- (158,290) -- (490,290) -- (490,284);
\node[anchor=center,inner sep=0pt,outer sep=0pt,text=black!93,font={\sffamily\fontsize{8}{9.6}\selectfont\bfseries}] at (324,308) {7 private};
\draw[fill=black!7,draw=black!40,line join=round,line width=0.599pt] (576,236) -- (593.321,226) -- (576,216) -- (558.679,226) -- cycle;
\draw[fill=cube1,draw=black!87,line join=round,line width=0.599pt] (576,236) -- (558.679,226) -- (558.679,246) -- (576,256) -- cycle;
\node[anchor=center,inner sep=0pt,outer sep=0pt,text=black!93,font={\sffamily\fontsize{8}{9.6}\selectfont}] at (567.34,241) {1};
\draw[fill=cube4,draw=black!87,line join=round,line width=0.599pt] (576,236) -- (576,256) -- (593.321,246) -- (593.321,226) -- cycle;
\node[anchor=center,inner sep=0pt,outer sep=0pt,text=black!93,font={\sffamily\fontsize{8}{9.6}\selectfont}] at (584.66,241) {4};
\node[anchor=center,inner sep=0pt,outer sep=0pt,text=black!93,font={\sffamily\fontsize{8}{9.6}\selectfont}] at (576,271) {14};
\draw[fill=black!7,draw=black!40,line join=round,line width=0.599pt] (624,236) -- (641.321,226) -- (624,216) -- (606.679,226) -- cycle;
\draw[fill=cube1,draw=black!87,line join=round,line width=0.599pt] (624,236) -- (606.679,226) -- (606.679,246) -- (624,256) -- cycle;
\node[anchor=center,inner sep=0pt,outer sep=0pt,text=black!93,font={\sffamily\fontsize{8}{9.6}\selectfont}] at (615.34,241) {1};
\draw[fill=cube4,draw=black!87,line join=round,line width=0.599pt] (624,236) -- (624,256) -- (641.321,246) -- (641.321,226) -- cycle;
\node[anchor=center,inner sep=0pt,outer sep=0pt,text=black!93,font={\sffamily\fontsize{8}{9.6}\selectfont}] at (632.66,241) {4};
\node[anchor=center,inner sep=0pt,outer sep=0pt,text=black!93,font={\sffamily\fontsize{8}{9.6}\selectfont}] at (624,271) {14};
\draw[fill=black!7,draw=black!40,line join=round,line width=0.599pt] (672,236) -- (689.321,226) -- (672,216) -- (654.679,226) -- cycle;
\draw[fill=cube1,draw=black!87,line join=round,line width=0.599pt] (672,236) -- (654.679,226) -- (654.679,246) -- (672,256) -- cycle;
\node[anchor=center,inner sep=0pt,outer sep=0pt,text=black!93,font={\sffamily\fontsize{8}{9.6}\selectfont}] at (663.34,241) {1};
\draw[fill=cube1,draw=black!87,line join=round,line width=0.599pt] (672,236) -- (672,256) -- (689.321,246) -- (689.321,226) -- cycle;
\node[anchor=center,inner sep=0pt,outer sep=0pt,text=black!93,font={\sffamily\fontsize{8}{9.6}\selectfont}] at (680.66,241) {1};
\node[anchor=center,inner sep=0pt,outer sep=0pt,text=black!93,font={\sffamily\fontsize{8}{9.6}\selectfont}] at (672,271) {11};
\draw[draw=black!87,line width=0.699pt] (554,284) -- (554,290) -- (694,290) -- (694,284);
\node[anchor=center,inner sep=0pt,outer sep=0pt,text=black!93,font={\sffamily\fontsize{8}{9.6}\selectfont\bfseries}] at (624,308) {one per axis};
\node[anchor=west,inner sep=0pt,outer sep=0pt,text=black!93,font={\sffamily\fontsize{8}{9.6}\selectfont}] at (716,216) {7 + \(\tau\) positions};
\node[anchor=west,inner sep=0pt,outer sep=0pt,text=black!93,font={\sffamily\fontsize{8}{9.6}\selectfont}] at (716,240) {need U or V;};
\node[anchor=west,inner sep=0pt,outer sep=0pt,text=black!93,font={\sffamily\fontsize{8}{9.6}\selectfont\bfseries}] at (716,264) {only \(\tau\)+6 copies.};
\draw[draw=black!20,line width=0.599pt] (16,320) -- (884,320);
\draw[fill=cube2,draw=black!87,line join=round,line width=0.599pt] (48,334) -- (72,334) -- (72,358) -- (48,358) -- cycle;
\node[anchor=center,inner sep=0pt,outer sep=0pt,text=black!93,font={\sffamily\fontsize{8}{9.6}\selectfont}] at (60,346) {2};
\draw[fill=cube4,draw=black!87,line join=round,line width=0.599pt] (24,358) -- (48,358) -- (48,382) -- (24,382) -- cycle;
\node[anchor=center,inner sep=0pt,outer sep=0pt,text=black!93,font={\sffamily\fontsize{8}{9.6}\selectfont}] at (36,370) {4};
\draw[fill=cube0,draw=black!87,line join=round,line width=0.599pt] (48,358) -- (72,358) -- (72,382) -- (48,382) -- cycle;
\node[anchor=center,inner sep=0pt,outer sep=0pt,text=black!93,font={\sffamily\fontsize{8}{9.6}\selectfont}] at (60,370) {0};
\draw[fill=cube1,draw=black!87,line join=round,line width=0.599pt] (72,358) -- (96,358) -- (96,382) -- (72,382) -- cycle;
\node[anchor=center,inner sep=0pt,outer sep=0pt,text=black!93,font={\sffamily\fontsize{8}{9.6}\selectfont}] at (84,370) {1};
\draw[fill=cube3,draw=black!87,line join=round,line width=0.599pt] (96,358) -- (120,358) -- (120,382) -- (96,382) -- cycle;
\node[anchor=center,inner sep=0pt,outer sep=0pt,text=black!93,font={\sffamily\fontsize{8}{9.6}\selectfont}] at (108,370) {3};
\draw[fill=cube3,draw=black!87,line join=round,line width=0.599pt] (48,382) -- (72,382) -- (72,406) -- (48,406) -- cycle;
\node[anchor=center,inner sep=0pt,outer sep=0pt,text=black!93,font={\sffamily\fontsize{8}{9.6}\selectfont}] at (60,394) {3};
\node[anchor=center,inner sep=0pt,outer sep=0pt,text=black!93,font={\sffamily\fontsize{9}{10.8}\selectfont\bfseries}] at (72,428) {(\(\tau\)+6) \(\times\) W};
\node[anchor=center,inner sep=0pt,outer sep=0pt,text=black!93,font={\sffamily\fontsize{8}{9.6}\selectfont}] at (72,450) {031423};
\draw[fill=cube0,draw=black!87,line join=round,line width=0.599pt] (180,376) -- (162.679,366) -- (162.679,386) -- (180,396) -- cycle;
\node[anchor=center,inner sep=0pt,outer sep=0pt,text=black!93,font={\sffamily\fontsize{8}{9.6}\selectfont}] at (171.34,381) {0};
\draw[fill=cube1,draw=black!87,line join=round,line width=0.599pt] (180,376) -- (180,396) -- (197.321,386) -- (197.321,366) -- cycle;
\node[anchor=center,inner sep=0pt,outer sep=0pt,text=black!93,font={\sffamily\fontsize{8}{9.6}\selectfont}] at (188.66,381) {1};
\draw[fill=cube2,draw=black!87,line join=round,line width=0.599pt] (180,376) -- (197.321,366) -- (180,356) -- (162.679,366) -- cycle;
\node[anchor=center,inner sep=0pt,outer sep=0pt,text=black!93,font={\sffamily\fontsize{8}{9.6}\selectfont}] at (180,366) {2};
\node[anchor=center,inner sep=0pt,outer sep=0pt,text=black!93,font={\sffamily\fontsize{8}{9.6}\selectfont}] at (180,411) {012};
\draw[fill=cube0,draw=black!87,line join=round,line width=0.599pt] (228,376) -- (210.679,366) -- (210.679,386) -- (228,396) -- cycle;
\node[anchor=center,inner sep=0pt,outer sep=0pt,text=black!93,font={\sffamily\fontsize{8}{9.6}\selectfont}] at (219.34,381) {0};
\draw[fill=cube3,draw=black!87,line join=round,line width=0.599pt] (228,376) -- (228,396) -- (245.321,386) -- (245.321,366) -- cycle;
\node[anchor=center,inner sep=0pt,outer sep=0pt,text=black!93,font={\sffamily\fontsize{8}{9.6}\selectfont}] at (236.66,381) {3};
\draw[fill=cube1,draw=black!87,line join=round,line width=0.599pt] (228,376) -- (245.321,366) -- (228,356) -- (210.679,366) -- cycle;
\node[anchor=center,inner sep=0pt,outer sep=0pt,text=black!93,font={\sffamily\fontsize{8}{9.6}\selectfont}] at (228,366) {1};
\node[anchor=center,inner sep=0pt,outer sep=0pt,text=black!93,font={\sffamily\fontsize{8}{9.6}\selectfont}] at (228,411) {031};
\draw[fill=cube0,draw=black!87,line join=round,line width=0.599pt] (276,376) -- (258.679,366) -- (258.679,386) -- (276,396) -- cycle;
\node[anchor=center,inner sep=0pt,outer sep=0pt,text=black!93,font={\sffamily\fontsize{8}{9.6}\selectfont}] at (267.34,381) {0};
\draw[fill=cube4,draw=black!87,line join=round,line width=0.599pt] (276,376) -- (276,396) -- (293.321,386) -- (293.321,366) -- cycle;
\node[anchor=center,inner sep=0pt,outer sep=0pt,text=black!93,font={\sffamily\fontsize{8}{9.6}\selectfont}] at (284.66,381) {4};
\draw[fill=cube3,draw=black!87,line join=round,line width=0.599pt] (276,376) -- (293.321,366) -- (276,356) -- (258.679,366) -- cycle;
\node[anchor=center,inner sep=0pt,outer sep=0pt,text=black!93,font={\sffamily\fontsize{8}{9.6}\selectfont}] at (276,366) {3};
\node[anchor=center,inner sep=0pt,outer sep=0pt,text=black!93,font={\sffamily\fontsize{8}{9.6}\selectfont}] at (276,411) {043};
\draw[fill=cube1,draw=black!87,line join=round,line width=0.599pt] (324,376) -- (306.679,366) -- (306.679,386) -- (324,396) -- cycle;
\node[anchor=center,inner sep=0pt,outer sep=0pt,text=black!93,font={\sffamily\fontsize{8}{9.6}\selectfont}] at (315.34,381) {1};
\draw[fill=cube3,draw=black!87,line join=round,line width=0.599pt] (324,376) -- (324,396) -- (341.321,386) -- (341.321,366) -- cycle;
\node[anchor=center,inner sep=0pt,outer sep=0pt,text=black!93,font={\sffamily\fontsize{8}{9.6}\selectfont}] at (332.66,381) {3};
\draw[fill=cube2,draw=black!87,line join=round,line width=0.599pt] (324,376) -- (341.321,366) -- (324,356) -- (306.679,366) -- cycle;
\node[anchor=center,inner sep=0pt,outer sep=0pt,text=black!93,font={\sffamily\fontsize{8}{9.6}\selectfont}] at (324,366) {2};
\node[anchor=center,inner sep=0pt,outer sep=0pt,text=black!93,font={\sffamily\fontsize{8}{9.6}\selectfont}] at (324,411) {132};
\draw[fill=cube1,draw=black!87,line join=round,line width=0.599pt] (372,376) -- (354.679,366) -- (354.679,386) -- (372,396) -- cycle;
\node[anchor=center,inner sep=0pt,outer sep=0pt,text=black!93,font={\sffamily\fontsize{8}{9.6}\selectfont}] at (363.34,381) {1};
\draw[fill=cube3,draw=black!87,line join=round,line width=0.599pt] (372,376) -- (372,396) -- (389.321,386) -- (389.321,366) -- cycle;
\node[anchor=center,inner sep=0pt,outer sep=0pt,text=black!93,font={\sffamily\fontsize{8}{9.6}\selectfont}] at (380.66,381) {3};
\draw[fill=cube3,draw=black!87,line join=round,line width=0.599pt] (372,376) -- (389.321,366) -- (372,356) -- (354.679,366) -- cycle;
\node[anchor=center,inner sep=0pt,outer sep=0pt,text=black!93,font={\sffamily\fontsize{8}{9.6}\selectfont}] at (372,366) {3};
\node[anchor=center,inner sep=0pt,outer sep=0pt,text=black!93,font={\sffamily\fontsize{8}{9.6}\selectfont}] at (372,411) {133};
\draw[fill=cube2,draw=black!87,line join=round,line width=0.599pt] (420,376) -- (402.679,366) -- (402.679,386) -- (420,396) -- cycle;
\node[anchor=center,inner sep=0pt,outer sep=0pt,text=black!93,font={\sffamily\fontsize{8}{9.6}\selectfont}] at (411.34,381) {2};
\draw[fill=cube3,draw=black!87,line join=round,line width=0.599pt] (420,376) -- (420,396) -- (437.321,386) -- (437.321,366) -- cycle;
\node[anchor=center,inner sep=0pt,outer sep=0pt,text=black!93,font={\sffamily\fontsize{8}{9.6}\selectfont}] at (428.66,381) {3};
\draw[fill=cube4,draw=black!87,line join=round,line width=0.599pt] (420,376) -- (437.321,366) -- (420,356) -- (402.679,366) -- cycle;
\node[anchor=center,inner sep=0pt,outer sep=0pt,text=black!93,font={\sffamily\fontsize{8}{9.6}\selectfont}] at (420,366) {4};
\node[anchor=center,inner sep=0pt,outer sep=0pt,text=black!93,font={\sffamily\fontsize{8}{9.6}\selectfont}] at (420,411) {234};
\draw[fill=cube3,draw=black!87,line join=round,line width=0.599pt] (468,376) -- (450.679,366) -- (450.679,386) -- (468,396) -- cycle;
\node[anchor=center,inner sep=0pt,outer sep=0pt,text=black!93,font={\sffamily\fontsize{8}{9.6}\selectfont}] at (459.34,381) {3};
\draw[fill=cube3,draw=black!87,line join=round,line width=0.599pt] (468,376) -- (468,396) -- (485.321,386) -- (485.321,366) -- cycle;
\node[anchor=center,inner sep=0pt,outer sep=0pt,text=black!93,font={\sffamily\fontsize{8}{9.6}\selectfont}] at (476.66,381) {3};
\draw[fill=cube4,draw=black!87,line join=round,line width=0.599pt] (468,376) -- (485.321,366) -- (468,356) -- (450.679,366) -- cycle;
\node[anchor=center,inner sep=0pt,outer sep=0pt,text=black!93,font={\sffamily\fontsize{8}{9.6}\selectfont}] at (468,366) {4};
\node[anchor=center,inner sep=0pt,outer sep=0pt,text=black!93,font={\sffamily\fontsize{8}{9.6}\selectfont}] at (468,411) {334};
\draw[fill=cube0!30,draw=black!40,line join=round,line width=0.599pt] (516,376) -- (498.679,366) -- (498.679,386) -- (516,396) -- cycle;
\node[anchor=center,inner sep=0pt,outer sep=0pt,text=black!60,font={\sffamily\fontsize{8}{9.6}\selectfont}] at (507.34,381) {0};
\draw[fill=cube2!30,draw=black!40,line join=round,line width=0.599pt] (516,376) -- (516,396) -- (533.321,386) -- (533.321,366) -- cycle;
\node[anchor=center,inner sep=0pt,outer sep=0pt,text=black!60,font={\sffamily\fontsize{8}{9.6}\selectfont}] at (524.66,381) {2};
\draw[fill=cube4!30,draw=black!40,line join=round,line width=0.599pt] (516,376) -- (533.321,366) -- (516,356) -- (498.679,366) -- cycle;
\node[anchor=center,inner sep=0pt,outer sep=0pt,text=black!60,font={\sffamily\fontsize{8}{9.6}\selectfont}] at (516,366) {4};
\node[anchor=center,inner sep=0pt,outer sep=0pt,text=black!60,font={\sffamily\fontsize{8}{9.6}\selectfont}] at (516,411) {024};
\draw[draw=black!87,line width=0.699pt] (158,424) -- (158,430) -- (490,430) -- (490,424);
\node[anchor=center,inner sep=0pt,outer sep=0pt,text=black!93,font={\sffamily\fontsize{8}{9.6}\selectfont\bfseries}] at (324,448) {7 private};
\draw[fill=black!7,draw=black!40,line join=round,line width=0.599pt] (576,376) -- (593.321,366) -- (576,356) -- (558.679,366) -- cycle;
\draw[fill=cube3,draw=black!87,line join=round,line width=0.599pt] (576,376) -- (558.679,366) -- (558.679,386) -- (576,396) -- cycle;
\node[anchor=center,inner sep=0pt,outer sep=0pt,text=black!93,font={\sffamily\fontsize{8}{9.6}\selectfont}] at (567.34,381) {3};
\draw[fill=cube4,draw=black!87,line join=round,line width=0.599pt] (576,376) -- (576,396) -- (593.321,386) -- (593.321,366) -- cycle;
\node[anchor=center,inner sep=0pt,outer sep=0pt,text=black!93,font={\sffamily\fontsize{8}{9.6}\selectfont}] at (584.66,381) {4};
\node[anchor=center,inner sep=0pt,outer sep=0pt,text=black!93,font={\sffamily\fontsize{8}{9.6}\selectfont}] at (576,411) {34};
\draw[fill=black!7,draw=black!40,line join=round,line width=0.599pt] (624,376) -- (641.321,366) -- (624,356) -- (606.679,366) -- cycle;
\draw[fill=cube3,draw=black!87,line join=round,line width=0.599pt] (624,376) -- (606.679,366) -- (606.679,386) -- (624,396) -- cycle;
\node[anchor=center,inner sep=0pt,outer sep=0pt,text=black!93,font={\sffamily\fontsize{8}{9.6}\selectfont}] at (615.34,381) {3};
\draw[fill=cube3,draw=black!87,line join=round,line width=0.599pt] (624,376) -- (624,396) -- (641.321,386) -- (641.321,366) -- cycle;
\node[anchor=center,inner sep=0pt,outer sep=0pt,text=black!93,font={\sffamily\fontsize{8}{9.6}\selectfont}] at (632.66,381) {3};
\node[anchor=center,inner sep=0pt,outer sep=0pt,text=black!93,font={\sffamily\fontsize{8}{9.6}\selectfont}] at (624,411) {33};
\draw[fill=black!7,draw=black!40,line join=round,line width=0.599pt] (672,376) -- (689.321,366) -- (672,356) -- (654.679,366) -- cycle;
\draw[fill=cube3,draw=black!87,line join=round,line width=0.599pt] (672,376) -- (654.679,366) -- (654.679,386) -- (672,396) -- cycle;
\node[anchor=center,inner sep=0pt,outer sep=0pt,text=black!93,font={\sffamily\fontsize{8}{9.6}\selectfont}] at (663.34,381) {3};
\draw[fill=cube4,draw=black!87,line join=round,line width=0.599pt] (672,376) -- (672,396) -- (689.321,386) -- (689.321,366) -- cycle;
\node[anchor=center,inner sep=0pt,outer sep=0pt,text=black!93,font={\sffamily\fontsize{8}{9.6}\selectfont}] at (680.66,381) {4};
\node[anchor=center,inner sep=0pt,outer sep=0pt,text=black!93,font={\sffamily\fontsize{8}{9.6}\selectfont}] at (672,411) {34};
\draw[draw=black!87,line width=0.699pt] (554,424) -- (554,430) -- (694,430) -- (694,424);
\node[anchor=center,inner sep=0pt,outer sep=0pt,text=black!93,font={\sffamily\fontsize{8}{9.6}\selectfont\bfseries}] at (624,448) {one per axis};
\node[anchor=west,inner sep=0pt,outer sep=0pt,text=black!93,font={\sffamily\fontsize{8}{9.6}\selectfont}] at (716,356) {7 + \(\tau\) positions};
\node[anchor=west,inner sep=0pt,outer sep=0pt,text=black!93,font={\sffamily\fontsize{8}{9.6}\selectfont}] at (716,380) {need W;};
\node[anchor=west,inner sep=0pt,outer sep=0pt,text=black!93,font={\sffamily\fontsize{8}{9.6}\selectfont\bfseries}] at (716,404) {only \(\tau\)+6 copies.};
\draw[draw=black!20,line width=0.599pt] (16,466) -- (884,466);
\end{tikzpicture}

\caption{The uniform obstruction $U+(\tau+5)V+(\tau+6)W$, where $\tau=a+b+c-6$.
Each net has $+x$ at its centre, and corner types are read
counterclockwise from outside. Only $U$, $V$ and $W$ can be targets.
Their pools are $U$ alone, $U$ or $V$, and $W$ alone; faded corner types
also occur on a cube outside the pool. For $V$ and $W$, one edge in each
direction has a private type, and these edges have $\tau$ noncorner positions
in total.}
\label{fig:uniform-obstruction}
\end{figure}

\begin{proposition}\label{cor:all-lower-bounds}
For every $a,b,c\ge2$,
\[
 F_5(a,b,c)\ge\max\{15,\;4(a+b+c)-16,\;2(a+b+c)+1\}.
\]
\end{proposition}
\begin{proof}
The three terms follow from the fourteen-cube corner obstruction (Lemma~\ref{lem:lower}),
the number of frame positions, and
Lemma~\ref{lem:thin-three-variety-obstruction}, respectively.
\end{proof}

\subsection{The smallest thin box}

\begin{theorem}[The $2\times2\times3$ box]\label{thm:223}
We have
\[
 F_5(2,2,3)=15.
\]
In particular, any fifteen five-colour cubes contain twelve that form a
$2\times2\times3$ box with monochromatic exterior faces using all five
colours.
\end{theorem}

\begin{proof}
The fourteen-cube corner obstruction (Lemma~\ref{lem:lower}) gives the lower bound: the eight
corner cubes of any box would themselves give a corner solution.

For the upper bound, take fifteen arbitrary source copies. Call a
bichromatic pair rare if at most seven sources possess it. Each rare
pair is missing from at least eight sources. The missing sets for any
two rare pairs intersect, since $8+8>15$. Lemma~\ref{lem:rare-pairs}
therefore gives at most two disjoint rare pairs.

Suppose first that there is no rare pair. The corner theorem supplies
a target $T$ and a corner matching. If $T$ is adjacent-repeat, choose
the long axis perpendicular to its monochromatic edge. If $T$ is
opposite-repeat, any choice of long axis is suitable. In either case
all four noncorner frame positions have bichromatic requirements, and
each requirement has at least eight source neighbours.

If there is a rare pair $e=\alpha\beta$, at least eight sources belong to
$\mathcal H_{\alpha\beta}$. Choose eight of them and any three additional copies.
Lemma~\ref{lem:eight-common-missing} gives a corner matching with target
$T\in\mathcal H_{\alpha\beta}$. Orient its $\alpha\beta$ axis along
one of the two short directions, and choose a perpendicular direction
as the long axis. The four noncorner edge requirements then each join
a pole colour, $\alpha$ or $\beta$, to a ring colour. None is $\alpha\beta$, and none can
be the other rare pair, if present, since that pair is disjoint from
$\alpha\beta$ and uses only ring colours. The target's monochromatic edge,
when present, is parallel to the $\alpha\beta$ axis and contributes no noncorner
position. Again all four edge requirements are bichromatic and have
at least eight source neighbours.

For either chosen target and orientation, every source fits at least
four of the twelve positions, by Lemma~\ref{lem:thin-compatibility}.
Thus each source misses at most eight positions.

Let $D$ be any subset of the twelve target positions. If $D$ contains
only corners, its Hall inequality follows from the corner matching.
Otherwise, if $|D|\le8$, one of its edge positions has at least eight
neighbours. If $|D|>8$, every source neighbours $D$, because no source
misses more than eight positions. Hall's theorem supplies a matching
of all twelve positions to distinct source copies. The final matching
may rearrange the preliminary corner placements.
\end{proof}

\subsection[Sharp thresholds for 2 by 3 by c frames]{Sharp thresholds for $2\times3\times c$ frames}

\begin{lemma}[Compatibility for $2\times3\times c$]\label{lem:23c-compatibility}
Let $M=c-2\ge2$ and put $d=2M+6$. Orient a five-colour target for
a $2\times3\times c$ frame so that its monochromatic edge, if any,
is parallel to the side of length two. An adjacent-repeat source misses
at most $d$ positions, and an opposite-repeat source misses at most
$d+1$ positions. If a source misses $d+1$ positions, it is the unique
opposite-repeat variety lacking the two bichromatic types appearing
among those missed edge positions.
\end{lemma}

\begin{proof}
The nonzero edge lengths are $1,M$. By Lemma~\ref{lem:cube-incidence},
for an opposite-repeat source,
at most six corners are incompatible. Its two missing bichromatic pairs
are disjoint, so at most three geometric edges are incompatible, with
at most two parallel to either nonzero direction. If their counts are
$k_1,k_M$, the defect is at most
\[
 6+k_1+Mk_M\le6+3+2(M-1)=2M+7=d+1.
\]
Equality requires six incompatible corners, two incompatible long
edges, and one incompatible unit-length edge. These three geometric
edges display exactly two bichromatic types: there are at least two
because any target type occurs at most twice, and at most two because
the source lacks only two types. They are disjoint and determine the
opposite-repeat source variety uniquely.

An adjacent-repeat source lacks only one bichromatic type. Against an
opposite-repeat target, at most six corners and edge positions of total
weight $2M$ are incompatible, giving defect at most $d$.
Against an adjacent-repeat target, an incompatible edge lies transverse
to the monochromatic edge, so Lemma~\ref{lem:transverse-corner}
gives a shared corner.
The absent type contributes weight at most $M+1$: a doubled type has
one edge in each nonzero direction. The defect is therefore at most
\[
 7+(M+1)=M+8\le2M+6=d.
\]
If no actual edge is incompatible, the defect is at most eight, which
is also at most $d$.
\end{proof}

\begin{lemma}[An opposite-repeat target along the short axis]
\label{lem:23c-opposite-short}
Let $c\ge3$. For a $2\times3\times c$ frame, orient an opposite-repeat target with
its repeated-colour axis along the side of length two. Every source
misses at most $2(c-2)+6$ positions.
\end{lemma}

\begin{proof}
Every actual edge type contains the target's repeated colour. By
Lemma~\ref{lem:cube-incidence}, a source
can lack at most one such bichromatic type, because its missing pairs
are disjoint. The absent type accounts for at most two parallel geometric
edges, of total weight at most $2(c-2)$. At least two corners are shared,
so at most six further positions are missed.
\end{proof}

\begin{theorem}[The remaining $2\times3\times c$ frames]
\label{thm:23c-sharp}
\[
 F_5(2,3,3)=17,\qquad
 F_5(2,3,c)=4c+4\quad(c\ge4).
\]
\end{theorem}

\begin{proof}
First let $c=3$ and take seventeen sources. Call a bichromatic pair
rare if at most eight sources possess it. Each rare pair is absent from
at least nine sources. By Lemma~\ref{lem:rare-pairs} there are
at most two, and they are disjoint. If no pair is rare, choose a corner target using $N_5=15$
and put its possible monochromatic edge along the side of length two.
Otherwise choose a rare pair $\alpha\beta$ and apply
Lemma~\ref{lem:eight-common-missing} to eight sources in $\mathcal H_{\alpha\beta}$
and any three others. Orient the target's $\alpha\beta$ axis along the side of
length two. All actual edge types join a pole colour to a ring colour,
so they avoid both possible rare pairs and are bichromatic.

Every source misses at most nine of the sixteen frame positions. Indeed,
if source or target is opposite-repeat, at most six corners and three
actual geometric edges are incompatible. If both are adjacent-repeat,
at most two actual geometric edges are incompatible; whenever one is
incompatible, Lemma~\ref{lem:transverse-corner} supplies a shared corner. If no edge
is incompatible, at most eight positions are missed.

The chosen target has a corner matching, and every edge requirement has
at least nine neighbours. Lemma~\ref{lem:partial-matching}, with the corners
as the matched subset, completes the frame. This proves $F_5(2,3,3)\le17$; the
collection $U+7V+8W$ in Lemma~\ref{lem:thin-three-variety-obstruction}
gives equality.

Now let $c\ge4$, put $M=c-2\ge2$, and set
\[
 d=2M+6,\qquad N=4M+12=2d.
\]
There are $N=4c+4$ frame positions. Take any $N$ source copies, and call
a bichromatic pair rare if at most $d-1$ copies possess it. Each rare
pair is absent from at least $d+1>N/2$ sources. Lemma~\ref{lem:rare-pairs}
gives at most two disjoint rare pairs.

If no pair is rare, use the corner theorem and orient the target's
possible monochromatic edge along the side of length two. Otherwise,
choose a rare pair $\alpha\beta$. At least $d+1\ge11$ sources belong to
$\mathcal H_{\alpha\beta}$. Lemma~\ref{lem:eight-common-missing} supplies a corner solution
whose target belongs to that family. Orient its $\alpha\beta$ axis along the side
of length two. All actual edge types join a pole colour to a ring colour,
and hence avoid both rare pairs; the monochromatic edge, if present,
has no noncorner positions.

In either case the chosen target has a corner matching and every actual
edge requirement has at least $d$ source neighbours. Hall's inequality
holds for corner-only sets, and for any set of at most $d$ positions
containing an edge. By Lemma~\ref{lem:23c-compatibility}, sets of more
than $d+1$ positions neighbour every source.

Suppose the frame has no matching. A deficient Hall set $D$ must therefore
have exactly $d+1$ positions. Every nonneighbour misses all of those
positions, so Lemma~\ref{lem:23c-compatibility} says that all nonneighbours
have one opposite-repeat variety $V$, determined by the two absent
bichromatic types. Since $D$ has at most $d$ neighbours among $N=2d$
copies, the inventory contains at least $d$ copies of $V$.

Retarget the frame to $V$ with its repeated-colour axis along the side
of length two. By Lemma~\ref{lem:23c-opposite-short}, every source misses
at most $d$ positions. The $d$ copies of $V$ are compatible with every
position, so they verify Hall's inequality for every set of size at most
$d$. Larger sets neighbour every source. This gives a frame after all.

Thus $N$ sources always suffice. The number of frame positions supplies
the matching lower bound.
\end{proof}
% Analytic resolution of the 2 x 2 x c family.
% Inputs: N_5=15; the eight-family-plus-three corner lemma; basic geometry.

\subsection[Sharp thresholds for 2 by 2 by c boxes]{Sharp thresholds for $2\times2\times c$ boxes}

Put $M=c-2$. The frame of a $2\times2\times c$ box consists of eight
corners and four parallel geometric edges, each with $M$ positions. Its
size is $L=4M+8=4c$. We call an orientation of an adjacent-repeat target
\emph{transverse} if its monochromatic edge is not parallel to the long
axis. For an opposite-repeat target, its \emph{axial} orientation has
the repeated faces perpendicular to the long axis.

\begin{lemma}[Directional compatibility]\label{lem:rod-compatibility}
Let $M\ge2$.
\begin{enumerate}
\item For a transverse adjacent-repeat target, every adjacent-repeat
source misses at most $M+7$ positions. Every opposite-repeat source
misses at most $2M+6$ positions. An opposite-repeat source missing more
than $M+6$ positions must be one of two varieties; their missing
long-edge types are complementary pairs.
\item For an axial opposite-repeat target $V$, every source misses at
most $d=2M+4=L/2$ positions. If $M\ge3$, a source different from $V$
can miss $d$ positions only if it is one of the other two opposite-repeat
varieties with the same repeated colour. These two varieties have
complementary missing long-edge types.
\end{enumerate}
\end{lemma}

\begin{proof}
For (i), write the target in face order $(+x,-x,+y,-y,+z,-z)$ as
\[
 (r,\alpha,r,\beta,\gamma,\delta),
\]
and take the long direction to be $x$. Its four long-edge types are
\[
 r\gamma,\quad r\delta,\quad \beta\gamma,\quad \beta\delta.
\]
An adjacent-repeat source lacks at most one of these. If it lacks one,
Lemma~\ref{lem:transverse-corner} gives a shared corner, so its defect
is at most $M+7$. If it lacks none, its defect is at most eight, also
bounded by $M+7$.

By Lemma~\ref{lem:cube-incidence}, an opposite-repeat source shares at
least two corners and lacks at most
two long-edge types, giving defect at most $2M+6$. To lack two of the
four displayed types, its two missing pairs must be one of the disjoint
matchings
\[
 \{r\gamma,\beta\delta\}\qquad\text{or}\qquad\{r\delta,\beta\gamma\}.
\]
Each matching determines exactly one opposite-repeat variety, with
repeated colour $\alpha$. The two varieties have complementary missing
long-edge types. Every other opposite-repeat source lacks at most one
long-edge type, and hence has defect at most $M+6$.

For (ii), let the repeated colour of $V$ be $r$, and put its singleton
opposite pairs on the two transverse axes. The four long-edge types are
the four edges of the resulting $K_{2,2}$ on the singleton colours.
The other two opposite-repeat varieties with repeated colour $r$ each
share four corners and two long-edge types with $V$. Their compatible
position count is $4+2M=d$, and their missing long-edge pairs are
complementary.

Every adjacent-repeat source lacks at most one long-edge type. An
opposite-repeat source with repeated colour different from $r$ also
lacks at most one: one of its two missing pairs contains $r$, whereas
no long-edge type does. By Lemma~\ref{lem:cube-incidence}, these sources share at least two
corners with $V$, so they fit at least $2+3M\ge d$ positions. If $M\ge3$, this last
inequality is strict. The target variety $V$ itself fits every position.
This proves both assertions.
\end{proof}

\begin{lemma}[Concentrated opposite-repeat copies]\label{lem:rod-concentration}
Let $M\ge3$ and $d=2M+4$. A collection of $2d=4M+8$ cubes containing
at least $d-1$ copies of one opposite-repeat variety forms a
$2\times2\times(M+2)$ box.
\end{lemma}

\begin{proof}
Let that variety be $V$ with multiplicity $v$. Use its axial target
orientation. Every source misses at most $d$ positions. If $v\ge d$,
Hall's inequality is immediate: sets of at most $d$ positions have the
$v$ copies of $V$ as neighbours, and larger sets neighbour every source.

Suppose $v=d-1$. The same reasoning handles all position sets except
sets $D$ of size exactly $d$. If such a set fails Hall, no non-$V$
copy neighbours it. Those copies must each miss $d$ positions, so
Lemma~\ref{lem:rod-compatibility}(ii) restricts them to the other two
opposite-repeat varieties with the same repeated colour. Since $d>8$,
$D$ contains a long-edge position. The two candidate varieties have
complementary missing long-edge types, so at most one can be absent
from the neighbourhood of $D$. All $2d-v=d+1$ non-$V$ copies are
therefore one variety $W$. Using $W$ as the axial target gives the
previous case, with at least $d$ universal copies, and completes the box.
\end{proof}

\begin{lemma}[A target avoiding scarce long-edge types]\label{lem:rod-target}
Suppose a collection has $N\ge15$ cubes. Call a bichromatic pair rare
when at most $\rho$ sources possess it, where $N-\rho\ge8$ and
$2(N-\rho)>N$. Then there is a corner-solvable target, transverse if
adjacent-repeat and axial if opposite-repeat, whose four long-edge
types are all non-rare.
\end{lemma}

\begin{proof}
Each rare pair is missing from more than half the collection. Thus any
two rare pairs are disjoint and there are at most two, by
Lemma~\ref{lem:rare-pairs}. If there is no
rare pair, use the corner theorem and choose either specified orientation.

Otherwise let $\alpha\beta$ be rare. At least eight sources belong to
$\mathcal H_{\alpha\beta}$. Lemma~\ref{lem:eight-common-missing} supplies a corner target
in this family. Its $\alpha\beta$ axis is a singleton opposite pair and can be
put along a short direction. If the target is opposite-repeat, choose
its repeated-colour axis as the long direction; otherwise choose either
axis perpendicular to the $\alpha\beta$ axis. Every long-edge type then joins
one of $\alpha,\beta$ to a ring colour. It is neither $\alpha\beta$ nor the other possible
rare pair, whose colours are both ring colours. These orientations are
axial or transverse as required.
\end{proof}

\begin{theorem}[Thin rods]\label{thm:sharp-rods-upper}
\[
 F_5(2,2,4)\le17,
 \qquad
 F_5(2,2,c)=4c\quad(c\ge5).
\]
\end{theorem}

\begin{proof}
First let $c=4$, so $M=2$ and $L=16$, and take seventeen source
copies. In Lemma~\ref{lem:rod-target} use $\rho=8$. The chosen target has
a corner matching, and every long-edge type has at least nine sources.
If the target is opposite-repeat, its defect is at most eight by
Lemma~\ref{lem:rod-compatibility}, and Hall's theorem completes it.

For an adjacent-repeat target, adjacent-repeat sources miss at most nine
positions and opposite-repeat sources at most ten. A failing Hall set
$D$ cannot consist only of corners. Since it contains an edge position,
it has at least nine neighbours; and if $|D|>10$, it neighbours every
source. Therefore a failure would have $|D|=10$ and exactly nine
neighbours. Its eight non-neighbouring source copies must all be
opposite-repeat. Each must miss two long-edge types, since otherwise
its defect is at most $M+6=8$. The two candidate varieties in
Lemma~\ref{lem:rod-compatibility}(i) have complementary missing
long-edge types. As $D$ contains an edge position, all eight copies
must be the same opposite-repeat variety $V$.

The axial target $V$ is solved by Hall directly: its eight copies
neighbour every nonempty set of at most eight positions, and larger
sets neighbour all seventeen sources, since each misses at most eight
positions. This proves the seventeen-cube upper bound.

Now let $M=c-2\ge3$, put $d=2M+4$, and take exactly $N=2d=L$
cubes. Use $\rho=d-1$ in Lemma~\ref{lem:rod-target}. Every long-edge type
of the chosen target has at least $d$ sources. An opposite-repeat
target is completed directly by Hall, since every source misses at
most $d$ positions.

For an adjacent-repeat target, adjacent-repeat sources miss at most
$M+7\le d$ positions; opposite-repeat sources miss at most
$2M+6=d+2$. Any failing Hall set $D$ therefore has
\[
 d+1\le |D|\le d+2.
\]
It has at most $d+1$ neighbours, leaving at least $d-1$ non-neighbouring
source copies. All are opposite-repeat. An opposite-repeat source
missing only one long-edge type has defect at most $M+6<d+1$, so
all these copies are among the two varieties in
Lemma~\ref{lem:rod-compatibility}(i). Since $D$ contains an edge,
all must be one variety $V$. Lemma~\ref{lem:rod-concentration} then
supplies a box from the original collection. Thus every collection of
$L=4c$ cubes works. The number of positions gives the matching lower
bound.
\end{proof}

\begin{proposition}[Complete rod threshold]\label{cor:sharp-rods}
Together with Theorem~\ref{thm:main}, Theorem~\ref{thm:223} and
the sixteen-cube obstruction of Lemma~\ref{lem:thin-three-variety-obstruction},
Theorem~\ref{thm:sharp-rods-upper} gives
\[
 F_5(2,2,c)=\max\{15,\,4c,\,2c+9\}\qquad(c\ge2).
\]
Equivalently, the values are $15,15,17$ for $c=2,3,4$, respectively,
and $4c$ for every $c\ge5$.
\end{proposition}
\begin{proof}
The cases $c=2,3$ are the established corner and smallest thin-box
thresholds. For $c=4$, Lemma~\ref{lem:thin-three-variety-obstruction} and the seventeen-cube
upper bound give equality. For $c\ge5$, the number of positions gives
$F_5(2,2,c)=4c$, as proved above.
\end{proof}
\subsection{Boxes whose two shortest sides sum to at least six}

Order the side lengths as $2\le a\le b\le c$ and put
\[
 m=a-2,\qquad t=b-2,\qquad M=c-2,\qquad \tau=m+t+M.
\]
The frame has $L=8+4\tau$ positions. In this section assume $m+t\ge2$.

\begin{lemma}[Weighted compatibility]\label{lem:box-weighted}
Let $2\le a\le b\le c$ with $a+b\ge6$, so that $m+t\ge2$ and $M\ge1$.
Every source misses at most $d=\tau+M+6$ positions of every target frame.
\end{lemma}
\begin{proof}
By Lemma~\ref{lem:cube-incidence}, if the target is opposite-repeat,
at most three geometric edges are
incompatible, with at most two in any direction. Their total weight is
at most $2M+t\le \tau+M$. At most six corners are incompatible, so the
claimed bound follows.

For an opposite-repeat source, take the adjacent-repeat target
\[
 (r,\alpha,r,\beta,\gamma,\delta),
\]
whose monochromatic edge is parallel to $z$.
A doubled bichromatic type, $r\gamma$ or $r\delta$, occupies one edge in each of
the $x,y$ directions. If such a type is missing, it and the monochromatic
edge contribute one edge in each direction; the other disjoint missing
type contributes at most one further edge. The missed weight is at most
$\tau+M$. If neither missing type is doubled, two missing single types
cannot have the same direction: the single types in each direction
share a colour. Adding the monochromatic edge therefore gives at most
three edges, with at most two parallel, again of weight at most
$2M+t\le \tau+M$. There are at most six incompatible corners.

For two adjacent-repeat varieties, at most three geometric edges are
incompatible. If there are at most two, the total defect is at most
\[
 2M+8\le \tau+M+6,
\]
because $m+t\ge2$. If there are three, they are the monochromatic edge
and a doubled missing bichromatic type, one edge in each direction.
Lemma~\ref{lem:transverse-corner}
gives a shared corner. The defect is at most
$\tau+7\le \tau+M+6$, since $M\ge1$.
\end{proof}

\begin{theorem}\label{thm:boxes-wide}
If $2\le a\le b\le c$ and $a+b\ge6$, then
\[
 F_5(a,b,c)=4(a+b+c)-16.
\]
\end{theorem}
\begin{proof}
Take $N=L=8+4\tau$ source copies and put $d=\tau+M+6$. Call a bichromatic
pair rare if at most $d-1$ copies possess it. Its missing family has
at least
\[
 N-d+1=3\tau-M+3
\]
copies, exceeding $N/2$ by $m+t-1\ge1$. Lemma~\ref{lem:rare-pairs}
again gives at most two disjoint rare pairs.

By Lemma~\ref{lem:partial-matching}, a target avoiding all rare pairs
has a full frame matching whenever its corners and monochromatic edge
positions have a matching: every other position has at least $d$
neighbours, and Lemma~\ref{lem:box-weighted} bounds each source's defect
by $d$.

\emph{No rare pair.} For each colour reserve up to $m$ sources having
that colour on adjacent faces, taking all such sources if fewer exist.
The reservations are disjoint. At least $N-5m\ge7m+8\ge15$ sources
remain when $m\ge1$; when $m=0$, none are reserved and $N\ge24$.
The corner theorem supplies a target on the unreserved sources.
Orient its possible monochromatic edge along the shortest direction.
If the target repeats colour $r$ adjacently, a compatible source with
adjacent $r$-faces survived the reservation, so the $r$-reserve has its
full size $m$. It fills the monochromatic edge. The Hall observation
completes the frame.

\emph{One rare pair $\alpha\beta$.} For each of the three ring colours,
reserve up to $m$ sources in the whole collection having that colour
on adjacent faces, again taking all such sources if fewer exist.
At most $3m$ family copies are removed, so $\mathcal H_{\alpha\beta}$ retains at least
\[
 (3\tau-M+3)-3m=3t+2M+3\ge8
\]
copies, and at least eleven total sources remain unreserved. Apply
Lemma~\ref{lem:eight-common-missing} to eight unreserved family copies
and any three further unreserved sources. The resulting corner target
lies in $\mathcal H_{\alpha\beta}$. Orient its $\alpha\beta$ axis along the shortest
direction. Any monochromatic edge is parallel to this axis and has its
full reserve of $m$ copies, by the same survivor argument. The target
avoids the only rare pair, so the Hall observation applies.

\emph{Two rare pairs $e,f$.} Their common missing variety $V$ is the
unique opposite-repeat variety with $e,f$ as singleton opposite pairs.
Its multiplicity satisfies
\[
 v\ge N-2(d-1)=2(m+t)-2\ge2.
\]
If $v\ge8$, use eight copies for the corners of $V$. If $4\le v\le7$,
the family $\mathcal H_e$ has at least eleven copies and hence at least
four not of variety $V$. Each shares at least four corners with $V$,
by the ring array. Place four such copies at distinct corners and fill
the other four with $V$. In either case, $V$ has no monochromatic edge
and avoids both rare pairs; the Hall observation applies.

It remains that $v=2$ or $3$, forcing $m+t=2$ and $N=2d$. If $n_e,n_f$
count the sources possessing the rare pairs, then
\[
 n_e,n_f\le d-1,\qquad
 v\ge N-n_e-n_f.
\]
Consequently,
\[
 n_e,n_f\ge d-2,\qquad n_f+v\ge d+1.
\]
Use exactly the reservation and the construction of Lemma~\ref{lem:eight-common-missing} from the
one-rare-pair case to obtain a corner-and-monochromatic-edge matching
for a target $T\in\mathcal H_e$, with its $e$ axis shortest.

Within this family, $V$ shares at least four target corners and can
miss at most two geometric edges, both parallel to the $e$ axis. For
a different opposite-repeat target these are its two $f$-edges; for
an adjacent-repeat target they are its $f$-edge and monochromatic edge.
Thus $V$ misses at most $2m+4\le6$ actual positions.

The only Hall sets not covered as before contain bichromatic edge
positions of type $f$ and no non-rare bichromatic edge. Such a set $D$
has at least $n_f$ neighbours. If $|D|>n_f\ge d-2\ge8$, every copy
of $V$ is also a neighbour, since it misses at most six positions.
These are distinct from the $n_f$ copies possessing $f$, so
\[
 |\Gamma(D)|\ge n_f+v\ge d+1.
\]
This handles $|D|\le d$; larger sets neighbour every source. Hall's
theorem completes the frame in the final case.

The number of frame positions gives the lower bound, proving equality.
All preliminary placements certify Hall inequalities; the final
matching may rearrange them.
\end{proof}

\subsection{Completion of the classification}
\begin{proof}[Proof of Theorem~\ref{thm:all-boxes}]
Order the sides as $2\le a\le b\le c$. When $a+b\ge6$,
Theorem~\ref{thm:boxes-wide} gives $F_5=L$. The remaining possibilities
are exactly $(a,b)=(2,2)$ and $(2,3)$, settled by
Proposition~\ref{cor:sharp-rods} and Theorem~\ref{thm:23c-sharp}.
Together with the corner theorem at $(2,2,2)$, these give $F_5=15$
for sums six and seven, $F_5=17$ for sum eight, and $F_5=L$
for every larger sum. Proposition~\ref{cor:all-lower-bounds} proves
sharpness and yields the stated maximum formula.

For the complete box, take $\max\{abc,F_5(a,b,c)\}$ sources and first
choose a frame. Every source has every colour, so any unused source can
fill any face-interior position after a suitable proper rotation. All
fully interior positions are unrestricted. There are enough unused
sources to fill the remaining $abc-L$ positions. Conversely, any complete
box supplies a frame and uses $abc$ sources. Hence
\[
 G_5(a,b,c)=\max\{abc,F_5(a,b,c)\}.
\]
Putting $x=a-2$, $y=b-2$, and $z=c-2$ gives
\[
 abc-L=xyz+2(xy+xz+yz)\ge0,
\]
which simplifies this identity to the claimed formula for $G_5$.

If $(a,b)=(2,2)$, this formula is $\max\{4c,15,2c+9\}$: it exceeds
the volume only for $c=2,3,4$. In every other sorted case, $b\ge3$;
then $abc\ge18\ge15$ and
\[
 abc-2(a+b+c)-1\ge 2bc-2b-2c-5\ge1.
\]
The first inequality follows by increasing $a$ from two, and the second
is minimised at $b=c=3$. There are no further exceptions.
\end{proof}

